\chapter{Duodenal Ulcer}
\index{Duodenal Ulcer}

\section*{Definition and Overview}
A duodenal ulcer is an open sore that forms in the lining of the duodenum, the first part of the small intestine, immediately beyond the stomach. It is a type of peptic ulcer, along with gastric ulcers, which form in the stomach itself. The most common causes are infection with the bacterium Helicobacter pylori and the use of non-steroidal anti-inflammatory medicines such as ibuprofen. Duodenal ulcers cause pain, particularly when the stomach is empty, and can lead to serious complications such as bleeding and perforation if untreated. Modern treatment is highly effective, and most ulcers heal completely.

\section*{Epidemiology}
Peptic ulcers, including duodenal ulcers, are common worldwide, although their prevalence has fallen over recent decades with better hygiene, lower rates of Helicobacter pylori infection in developed countries, and the widespread use of acid-suppressing medicines. Duodenal ulcers are more common than gastric ulcers and have historically affected men more than women, though the gap has narrowed. Complications such as bleeding are a significant cause of emergency hospital admission, particularly in older people taking anti-inflammatory medicines.

\section*{Aetiology and Pathophysiology}
The duodenal lining is normally protected by a layer of mucus and bicarbonate, balanced against the acid produced by the stomach. An ulcer forms when this balance is disturbed.
\begin{itemize}
    \item \textbf{Helicobacter pylori infection:} the most common cause; this bacterium colonises the stomach lining, damages the protective mucus layer, and increases acid production
    \item \textbf{NSAIDs:} medicines such as ibuprofen, naproxen, and aspirin reduce the production of prostaglandins that protect the lining, increasing ulcer risk
    \item \textbf{Acid and pepsin:} the digestive juices erode the unprotected lining to form an ulcer
    \item \textbf{Smoking:} increases acid production and impairs healing
    \item \textbf{Alcohol:} irritates the lining and may slow healing
    \item \textbf{Stress and severe illness:} critical illness can cause stress ulcers, although ordinary psychological stress is a minor factor
    \item \textbf{Zollinger-Ellison syndrome:} a rare tumour causing excessive acid production
\end{itemize}

\section*{Clinical Features}
Duodenal ulcers cause characteristic symptoms, though some people have few or none.
\begin{itemize}
    \item A burning or gnawing pain in the upper abdomen, usually between meals and at night when the stomach is empty
    \item Pain relieved by eating or taking antacids
    \item Belching, bloating, and a feeling of fullness
    \item Nausea and sometimes vomiting
    \item Indigestion and heartburn
    \item Silent ulcers may present only when a complication such as bleeding occurs
    \item Warning signs of complications include black, tarry stools, vomiting blood, and sudden severe abdominal pain
\end{itemize}

\section*{Differential Diagnosis}
Upper abdominal pain has many causes that must be distinguished from a duodenal ulcer.
\begin{itemize}
    \item Gastric ulcer, which tends to cause pain during or soon after meals
    \item Gastro-oesophageal reflux disease and heartburn
    \item Gallstones and biliary colic
    \item Pancreatitis and pancreatic disease
    \item Irritable bowel syndrome and functional dyspepsia
    \item Angina, which can cause upper abdominal or chest discomfort
    \item Stomach cancer, which is rare but must be considered with weight loss or bleeding
\end{itemize}

\section*{When to Seek Medical Care}
Anyone with persistent upper abdominal pain, indigestion that does not settle with simple measures, or a family history of ulcer or stomach cancer should see a GP. People who take NSAIDs or aspirin regularly and develop indigestion should also seek advice, as should those with black stools or blood in vomit.

\textbf{Contact your local emergency services for:}
\begin{itemize}
    \item Vomiting blood or material that looks like coffee grounds
    \item Black, tarry, or bloody stools
    \item Sudden, severe, tearing abdominal pain, which may indicate a perforated ulcer
    \item Faintness, dizziness, or collapse, which may indicate significant bleeding
    \item Abdominal pain with a rigid, board-like abdomen and fever
\end{itemize}

\section*{Diagnosis and Investigations}
Investigation aims to confirm the ulcer and identify its cause.
\begin{itemize}
    \item \textbf{Endoscopy (gastroscopy):} a camera passed through the mouth into the stomach and duodenum directly visualises the ulcer and allows biopsy
    \item \textbf{H. pylori testing:} urea breath test, stool antigen test, or biopsy taken at endoscopy
    \item \textbf{Blood tests:} full blood count to check for anaemia from bleeding
    \item \textbf{Imaging:} CT scans in emergencies such as suspected perforation
    \item \textbf{Medication review:} assessment of NSAID and aspirin use
\end{itemize}

\section*{Management}
Treatment of duodenal ulcers is highly successful.
\begin{itemize}
    \item \textbf{Acid-suppressing medicines:} proton pump inhibitors such as omeprazole, or H2 blockers such as ranitidine, heal most ulcers within four to eight weeks
    \item \textbf{H. pylori eradication:} a two-week course of combination therapy, usually a proton pump inhibitor with two antibiotics, cures the infection and prevents recurrence
    \item \textbf{Stopping NSAIDs:} where possible, alternative pain relief is used, and protective medicines are prescribed if NSAIDs must continue
    \item \textbf{Lifestyle measures:} stopping smoking, limiting alcohol, and eating regular meals
    \item \textbf{Follow-up endoscopy:} in selected cases to confirm healing, particularly for gastric ulcers
    \item \textbf{Emergency treatment:} for bleeding ulcers, endoscopy with injection or clipping; for perforation, emergency surgery
\end{itemize}

\section*{Complications}
\begin{itemize}
    \item Bleeding, which may be slow and cause anaemia or sudden and life-threatening
    \item Perforation, when the ulcer erodes through the intestinal wall, causing peritonitis and requiring emergency surgery
    \item Obstruction, when scarring narrows the outlet of the stomach
    \item Recurrence, which is now uncommon when H. pylori is eradicated and NSAIDs are managed
\end{itemize}

\section*{Prognosis and Outcomes}
The outlook for duodenal ulcer is excellent with modern treatment. Ulcers heal within weeks, and eradicating H. pylori reduces the recurrence rate from around 70 per cent to below 10 per cent. The main risks relate to bleeding and perforation, which are serious emergencies, and to long-term NSAID use, which requires careful management. Most people who have had a duodenal ulcer can live normally once the cause is treated.

\section*{Prevention and Self-Care}
\begin{itemize}
    \item Use NSAIDs and aspirin at the lowest effective dose and shortest duration, and only as directed
    \item Take protective acid-suppressing medicines if you need long-term NSAID treatment
    \item Eat regular meals and avoid heavy, late-night meals
    \item Stop smoking and limit alcohol
    \item See a doctor for persistent indigestion rather than self-treating for months
    \item Seek urgent care for any signs of bleeding
\end{itemize}

\section*{Sources and Further Reading}
The following reputable sources provide further information for healthcare students and patients seeking reliable, current advice about duodenal ulcers.
\begin{itemize}
    \item NHS. \textit{Stomach ulcer: symptoms, causes, and treatment.} https://www.nhs.uk/conditions/stomach-ulcer/
    \item Mayo Clinic. \textit{Peptic ulcer: symptoms and causes.} https://www.mayoclinic.org/diseases-conditions/peptic-ulcer/
    \item MedlinePlus. \textit{Peptic ulcer.} https://medlineplus.gov/pepticulcer.html
    \item MSD Manual. \textit{Peptic ulcer disease: professional overview.} https://www.msdmanuals.com/
    \item Healthdirect Australia. \textit{Stomach and duodenal ulcers.} https://www.healthdirect.gov.au/stomach-ulcer
    \item Australian Gastroenterology Week resources. \textit{Ulcer information.} https://www.gesa.org.au/
\end{itemize}
